\scanpage{41}
Soit $A$ une VA sur $k$ alg. clos, d'où module $\mathbb{T}_\ell(A)$ \deleted{$\unsure{\otimes}$} sur la pro-alg. de Lie \deleted{$\unsure{\mathfrak{g}_\ell}$} $\mathfrak{g}(k)=\mathfrak{g}$.

Soit $K$ une clôture alg. de $\mathbb{Q}_\ell$, posant
\[
\mathbb{T}_K(A)=\mathbb{T}_\ell(A)
\]
de sorte que $\mathbb{T}_K(A)$ est un module sur $\mathfrak{g}\otimes_{\mathbb{Q}_\ell}K=\mathfrak{g}_K$. Supposons

a) La \deleted{\unsure{les}} conj de Tate vérifiée pr $H^2(A,\mathbb{Q}_\ell(1))$
\[
H^2(A\deleted{\unsure{\times}},\mathbb{Q}_\ell(1))^{\mathrm{alg}}
=H^2(A,\mathbb{Q}_\ell(1))^{\mathfrak{g}}
\]

b) $\mathfrak{g}_K$ opérant sur \deleted{$\unsure{\mathfrak{g}}$ est} $\mathbb{T}_K(A)$ est ``\unsure{pseudoalgébrique}''

[i.e. les valeurs propres des frobénius \deleted{\unsure{sont}} de norme $q$ sont $\varepsilon q^{-1/2}$, $\varepsilon$ une racine de l'unité\ldots]

\margin{\unsure{Pr. Tits}, le premier cas est un cas particulier du second}

\uline{ou} $\mathbb{T}_K(A)$ se décompose en deux morceaux $V_1\times V_2$, tels que chacun soit propre pr ``presque tous'' les Frobénius, avec des valeurs propres différentes pr presque tous les Frobénius.

[Ces conditions sont remplies si $A$ est une \deleted{produit}\inserted{puissance} \margin{$C^n$} d'une \deleted{de} courbe elliptique \inserted{$C$ définissable} sur un corps fini $\mathbb{F}_q$, comme on voit en notant que a) \deleted{\unsure{et b)} \unsure{triviale|triviales} \unsure{se réduit pr passage à une puiss.|se déduit pr passage à une puiss.} \unsure{adéquate, et pr lequel b) est}} a été vérifié pr Mumford. Ceci est connu.

\mspage{40}
